\uuid{KTKi}
\titre{Concave/convexe}
\auteur{bodin ; hanani ; mzari}
\organisation{exo7}
\licence{CC BY-NC-SA 4.0}
\datecreate{2020-01-20}
\niveau{L1}
\module{Analyse}
\chapitre{Dérivabilité des fonctions réelles}
\sousChapitre{Autre}
\difficulte{3}
\theme{Dérivées, concave, convexe}
\exo7qcmid{lille-1.derivees.028}

\contenu{
\begin{question}
\qclassification{124.00}
\qtags{motcle=concave/convexe}
Soit $f(x)=\mathrm{e}^x$. Quelles sont les bonnes réponses ?
\begin{answers}
    \bad{$f''(x)$ s'annule au moins une fois sur $\Rr$.}
    \good{$f$ est convexe sur $\Rr$.}
    \bad{$f$ est concave sur $\Rr$.}
    \good{$\forall t\in [0,1]$ et $\forall x,y\in \Rr^{+*}$, on a : $t\ln x+(1-t)\ln y\leq \ln \left[tx+(1-t)y\right]$.}
\end{answers}
\begin{explanations}
On a $f''(x)=\mathrm{e}^x>0$. Donc $f$ est convexe sur $\Rr$. Ainsi, par définition,
$$\forall t\in [0,1]\; \forall a,b\in \Rr,\; f\left[ta+(1-t)b\right]\leq tf(a)+(1-t)f(b).$$
En prenant $a=\ln x$ et $b=\ln y$, avec $x,y>0$, on aura
$$\mathrm{e}^{t\ln x+(1-t)\ln y}\leq tx+(1-t)y.$$
Il suffit de composer par ln, qui est strictement croissante, pour avoir
$$t\ln x+(1-t)\ln y\leq \ln \left[tx+(1-t)y\right].$$
\end{explanations}
\end{question}
}
